%% Career-Ops LaTeX CV Template
%% Style: ATS-optimized single-column, minimal color accent, single-page friendly.
%% Build: pdflatex {file}.tex (or tectonic {file}.tex)

\documentclass[letterpaper,10pt]{article}

\usepackage[margin=0.55in]{geometry}
\usepackage[T1]{fontenc}
\usepackage[utf8]{inputenc}
\usepackage{lmodern}
\usepackage{microtype}
\usepackage{enumitem}
\usepackage{titlesec}
\usepackage{hyperref}
\usepackage{xcolor}
\usepackage{parskip}

\definecolor{accent}{HTML}{1F3A5F}

\hypersetup{
  colorlinks=true,
  urlcolor=accent,
  linkcolor=accent,
  pdftitle={ Ahnaf An Nafee - Resume },
  pdfauthor={ Ahnaf An Nafee },
  pdfsubject={Resume},
  pdfkeywords={ Kubernetes, OpenShift, AWS, Python, C++, PyTorch, GPU clusters, distributed training, inference, LLMs, MLOps, observability, RBAC, container orchestration, CI/CD, autoscaling, Terraform, GitOps, SRE, AI Infrastructure }
}

% Ensure generated PDF is machine readable / ATS parseable
\input{glyphtounicode}
\pdfgentounicode=1

\pagestyle{empty}
\setlength{\parindent}{0pt}
\raggedbottom
\raggedright

% --- Section formatting ------------------------------------------------------
% Note: section titles render in mixed case (no \MakeUppercase) for maximum
% ATS compatibility. Older parsers normalize headers like "Work Experience"
% and "Education" more reliably than uppercase variants.
\titleformat{\section}
  {\color{accent}\large\bfseries}
  {}{0em}{}[{\color{accent}\vspace{-6pt}\hrulefill\vspace{-4pt}}]
\titlespacing*{\section}{0pt}{8pt}{4pt}

\titleformat{\subsection}
  {\normalsize\bfseries}
  {}{0em}{}
\titlespacing*{\subsection}{0pt}{4pt}{2pt}

% --- Custom commands ---------------------------------------------------------
\newcommand{\role}[4]{%
  \textbf{#1} \hfill \textit{#3}\\[-1pt]
  \textit{#2} \hfill \textit{#4}\par\vspace{2pt}
}

\newcommand{\education}[4]{%
  \textbf{#1} \hfill \textit{#3}\\[-1pt]
  \textit{#2} \hfill \textit{#4}\par\vspace{2pt}
}

\newlist{tightitemize}{itemize}{1}
\setlist[tightitemize]{leftmargin=1.3em, topsep=2pt, itemsep=1pt, parsep=0pt, label=\textbullet}

% =============================================================================
\begin{document}

% --- Header ------------------------------------------------------------------
\begin{center}
  {\LARGE\bfseries\color{accent} Ahnaf An Nafee }\\[3pt]
  {\normalsize Software Engineer \textbar{} AI Researcher \textbar{} DevOps \& Cloud Infrastructure }\\[4pt]
  \small
  \href{mailto:ahnafnafee@gmail.com}{ahnafnafee@gmail.com} \textbar{}
  +1-540-252-8738 \textbar{}
  Fairfax, VA\\
  \href{https://linkedin.com/in/ahnafnafee}{linkedin.com/in/ahnafnafee} \textbar{}
  \href{https://github.com/ahnafnafee}{github.com/ahnafnafee} \textbar{}
  \href{https://www.ahnafnafee.dev}{ahnafnafee.dev} \textbar{} \href{https://scholar.google.com/citations?user=u15DO0cAAAAJ}{Google Scholar}
\end{center}

% --- Summary -----------------------------------------------------------------
\section*{Summary}
Computer Science PhD candidate and infrastructure engineer with 4+ years building production Kubernetes/OpenShift platforms at enterprise scale, AWS cloud-native architectures, and CI/CD automation pipelines. Led cross-team standardization of K8s deployment dictionaries across 10+ service teams, drove an AWS migration that cut cloud spend by 80\% via right-sizing and autoscaling, and automated release workflows that eliminated 85\% of manual toil. Active research at George Mason's DCXR Lab spans PyTorch, diffusion models, LLMs, and distributed training---bridging applied AI/ML literacy with deep container-orchestration, observability, and RBAC fundamentals. Targeting AI infrastructure and ML platform roles where Kubernetes scheduler depth, GPU-cluster operations, model serving, and MLOps observability intersect.

% --- Skills ------------------------------------------------------------------
\section*{Technical Skills}
\textbf{Cloud \& Infrastructure:} AWS, Kubernetes, OpenShift, Docker, Terraform, Helm, GitOps, CI/CD, GitHub Actions, Jenkins, Gradle, Maven, Observability, Autoscaling\\
\textbf{AI / ML:} PyTorch, TensorFlow, LLMs, Diffusion Models, Generative AI, Deep Learning, Computer Vision, Reinforcement Learning, MLOps\\
\textbf{Languages:} Python, C++, Go, Java, TypeScript, JavaScript, SQL, GLSL, Bash\\
\textbf{Backend \& Systems:} Distributed Systems, Microservices, REST, gRPC, WebSockets, OAuth, RBAC, System Design, High Availability\\
\textbf{Graphics \& XR:} 3D Computer Graphics, Real-Time Rendering, Shader Programming, Unity, Unreal Engine, WebGL, Procedural Generation\\
\textbf{Frontend:} React, Node.js, Next.js, Redux, HTML5, CSS3\\
\textbf{Methods:} Agile, Scrum, Test-Driven Development, Code Review, Technical Leadership, Cross-Functional Collaboration

% --- Experience --------------------------------------------------------------
\section*{Work Experience}
\role{DevOps Engineer}{Paychex}{Feb 2023 -- Aug 2025}{Rochester, NY}
\begin{tightitemize}
  \item Architected and led cross-team standardization of Kubernetes/OpenShift deployment dictionaries across 10+ service teams, eliminating configuration conflicts and accelerating developer velocity at enterprise scale
  \item Designed and rolled out Role-Based Access Control (RBAC) policies in OpenShift in partnership with security and infrastructure teams, hardening the platform against privilege-escalation risks
  \item Built a Gradle plugin for container image certification that streamlined CI/CD workflows, optimized container registries and multi-stage builds, and lowered operational expenses by \textbf{10\% annually}
  \item Automated end-to-end TLS certificate lifecycle management, eliminating manual renewals and reducing certificate-expiration-related downtime by \textbf{100\%}
  \item Championed infrastructure-as-code and observability best practices; partnered with SREs to improve deployment reliability, monitoring/alerting coverage, and mean time to recovery (MTTR)
\end{tightitemize}

\role{Chief Technology Officer}{Dynasty 11 Studios}{Jun 2022 -- Feb 2023}{Wayne, PA}
\begin{tightitemize}
  \item Spearheaded migration of monolithic backend to a microservices architecture on AWS, cutting application load and cloud spend by \textbf{80\%} through right-sizing, autoscaling strategies, and serverless adoption
  \item Introduced CI/CD with GitHub Actions and Maven, automating builds, tests, and deployments; reduced manual release effort by \textbf{85\%} and eliminated deploy-day toil
  \item Owned end-to-end technical strategy, architecture, and execution as CTO; led engineering across backend, frontend, and cloud teams shipping a consumer-facing matchmaking platform
  \item Recruited, mentored, and led multiple engineering pods through full microservice design, implementation, and production rollout
  \item Drove weekly stakeholder syncs, aligning product, design, and engineering on roadmap, scope, and delivery milestones
\end{tightitemize}

\role{Software Engineer}{Dynasty 11 Studios}{Sep 2021 -- Jun 2022}{Wayne, PA}
\begin{tightitemize}
  \item Engineered a real-time chat subsystem using STOMP over WebSockets in Java, supporting thousands of concurrent users with low-latency delivery
  \item Integrated 20+ REST endpoints across third-party and OAuth-based identity providers, enabling secure, scalable authentication and data exchange
  \item Optimized client-side state management with Redux, cutting chat-service performance bottlenecks by \textbf{80\%}
  \item Designed data pipelines and APIs powering the matchmaking recommendation engine serving thousands of active users
\end{tightitemize}

\role{Graduate Teaching Assistant}{George Mason University, DCXR Lab}{Aug 2025 -- Present}{Fairfax, VA}
\begin{tightitemize}
  \item Mentor 60+ graduate and undergraduate students in algorithms, systems, and applied machine learning, driving measurable improvements in lab completion and conceptual mastery
  \item Design coursework, lab exercises, and evaluation rubrics aligned with modern software engineering best practices
\end{tightitemize}

\role{Technical Programmer}{PHL Collective}{Mar 2021 -- Sep 2021}{Philadelphia, PA}
\begin{tightitemize}
  \item Shipped gameplay systems and parameterized shaders for \textit{DC's Justice League: Cosmic Chaos} (multi-platform console/PC release)
  \item Built modular game-manager scripts that streamlined designer workflows and accelerated feature iteration
  \item Partnered with technical artists to create customizable, performance-tuned shaders optimized for runtime rendering
\end{tightitemize}

% --- Education ---------------------------------------------------------------
\section*{Education}
\education{Ph.D. in Computer Science}{George Mason University --- DCXR Lab}{}{Fairfax, VA}
\begin{tightitemize}
  \item Research focus: generative AI for 3D content creation, machine-learning-driven graphics pipelines (UV mapping, non-photorealistic rendering), and human--computer interaction in immersive XR environments
\end{tightitemize}

\education{B.S. in Computer Science}{Drexel University}{}{Philadelphia, PA}
\begin{tightitemize}
  \item Honors: Dean's List; Founder's Scholarship
  \item Relevant coursework: Algorithms, Operating Systems, Distributed Computing, Computer Graphics, Machine Learning, Software Engineering
\end{tightitemize}

% --- Projects ----------------------------------------------------------------
% Section header is "Selected Projects" (not "Research & Selected Projects")
% so ATS classify the items below as Projects. Research focus stays in the
% bullet text itself.
\section*{Selected Projects}
\begin{tightitemize}
\item \textbf{AI-Driven 3D Content Generation} --- Researching generative models (diffusion, neural fields) for automated, controllable 3D asset creation in production-grade graphics pipelines, leveraging PyTorch and distributed training workflows
\item \textbf{ML for Graphics Pipelines} --- Exploring learning-based approaches to UV mapping, stylization, and non-photorealistic rendering (NPR) for real-time applications
\item \textbf{Portfolio \& Publications} --- Additional projects, technical artwork, and engineering work available at \href{https://www.ahnafnafee.dev/portfolio}{ahnafnafee.dev/portfolio}
\end{tightitemize}

\end{document}
